\chapter{从任意输入到平方核心，再到一个有限乘法群}\label{ch:entry}
\section{先把所有同余见证集中到一个单点}
上一讲假设已经有 $(t^2,t,1)$ 核心。现在说明任何合法输入都能进入它。这部分已经对所有 $t\ge2$ 证明，而且不依赖后面的单位群结论。\src{S14}\src{S15}

记 $n=t^2+t+1$，$d=\omega(n)$ 为不同素因子个数。因为
\[2^{t+1}>t^2+t+1,
\]
必有 $d\le t$。固定一个位置，为每个 $\ell\mid n$ 再保护一个与它模 $\ell$ 不同的位置。总共至多保护 $t+1$ 个位置，余下足够平均一次，得到
\[(a^{t^2},c_1,\ldots,c_{t+1}).
\]
令 $\delta_i=c_i-a$。每个危险素数处，至少有一个 $\delta_i$ 非零。问题是：能不能把这些可能分散在不同单点上的非零信息集中起来？

前篇的块与单点交换只需一个大块。平均 $t^2-1$ 份 $a$ 和 $c_i$，将剩下的一份旧 $a$ 登记为新单点。差分作用为
\[A_i=I+(\lambda\mathbf1-e_i)e_i^T,\qquad\lambda=-1/t^2.
\]
三个不同角色 $i,j,h$ 给出交换子
\[[A_h^{-1}A_jA_h,A_i]
=I+\frac{t^2-1}{t^2}(e_h-e_j)e_i^T.
\]
只为调整模 $\rad(n)$ 的信息时，形式逆可由有限群中的正整数幂替代。得到的实际数值未必是有理逆的数值，但余数恰好相同，足以完成这一入口任务。

这个交换子能给 $\delta_j$ 加上 $\delta_i$ 的倍数，唯一可能失效的素数由
\[\gcd(t^2-1,n)=\gcd(t+2,3)
\]
控制。因此至多只有三失效。若三失效，先把一个模三非零的角色选作目标 $j$。其他危险素数逐个处理：在当前素数处选非零的 $\delta_i$，用CRT选择交换子重复次数，让这一次只改变当前素数处的目标余数。第三个单点提供缓冲。最后同一个 $\delta_j$ 在全部危险素数处非零。

留下目标单点 $w=c_j$ 和 $t$ 份旧 $a$，平均其余 $t^2$ 项。所得 $u,v$ 满足
\[v=a,\qquad u=(-ta-w)/t^2,\qquad
u-v\equiv(a-w)/t^2\pmod\ell.
\]
右边在每个 $\ell\mid n$ 处非零，于是到达合法平方核心。

这里的复用值得注意：旧证明在 $n>2k$ 中使用这一交换子，是因为后面还要建第二个大块。交换子本身只需要三个单点。把它的真实前提分离出来，就能用于 $n=t^2+t+1<2t^2$ 的新范围。

\section{剩下的是哪个有限群}
定义
\begin{equation}\label{eq:Ht}
H_t=\langle-1,\ d\bmod n:\ 1\le d\le t,\ (d,n)=1\rangle
\subseteq(\Z/n\Z)^\times.
\end{equation}
这个群只记录上一讲已经可以正向、可逆实现的倍率。局部化环 $R_t$ 的单位在模 $n$ 下的像恰好是 $H_t$。

完整上剪切和 $n$ 倍下剪切，结合前篇的同余工具，给出
\[\langle U(R_t),L(nR_t)\rangle=\Gamma_1(n,R_t),
\]
其中 $\Gamma_1$ 中矩阵的左下项为零模 $n$，两个对角项为一模 $n$。它能细调当前列，却不能任意改变第二坐标的模 $n$ 余数。改变这个余数正是倍率的工作。

\begin{proposition}[平方族的充分算术接口]\label{prop:Ht}
若 $H_t=(\Z/n\Z)^\times$，则 $t^2$ 平均在 $n=t^2+t+1$ 个位置上的完整G判据成立。
\end{proposition}
\begin{proof}
先用上一节进入合法核心，将参数取本原整数代表 $(x,y)$。由 $\gcd(y,n)=1$，可求
\[\alpha x+\beta y=1,\qquad n\mid\alpha.
\]
令 $F=\left(\begin{smallmatrix}y&-x\\\alpha&\beta\end{smallmatrix}\right)$，则 $F(x,y)^T=(0,1)^T$。

由 $H_t$ 满群，可以找到 $R_t$ 的单位 $\epsilon\equiv y\pmod n$。矩阵
\[D_\epsilon F\in\Gamma_1(n,R_t),\qquad
D_\epsilon=\diag(\epsilon^{-1},\epsilon)
\]
的每一项同余条件都成立：左上项为 $y/\epsilon\equiv1$，右下项由行列式一也为一，左下项仍被 $n$ 整除。$D_\epsilon$ 射影等价于已可执行的 $\Delta(\epsilon^2)$，故 $F$ 可以正向实现。

输出方向 $(0,1)$ 对应 $(0^{t^2},(-1)^t,t)$。选全部 $t$ 份负一、单点 $t$ 及 $t^2-t-1$ 个零，做一次平均即全零。
\end{proof}

这个命题把“还差一点算术”变成了可独立研究的有限群问题。其正向实现、任意输入入口和最终操作都已明确。

\section{为什么换更多同形终点并没有绕过它}
如果 $H_t$ 不满，能否换一个终端方向？先看所有在本核心内一次选出零和组的方向。

选 $b$ 份 $v$，不选单点时，零和方程为 $t^2x-by=0$；本原方向的第二坐标是
\[y=\frac{t^2}{\gcd(b,t^2)},\qquad0\le b\le t.
\]
若选单点，方程为 $-(t+1)x+(t-b)y=0$，第二坐标是
\[y=\frac{t+1}{\gcd(t-b,t+1)}.
\]
两种 $y$ 的模 $n$ 余数都属于 $H_t$。第一种只用不超过 $t$ 的因子；第二种还用到
\[t+1\equiv-t^2\pmod n.
\]
边界 $b=0$ 或 $b=t$ 直接给 $y=1$。

所以继续枚举这种一步零和方向，不能补齐 $H_t$ 外的尺度类。这个限制针对当前核心和当前倍率菜单。它没有排除离开核心再返回、用更长终端词，或通过状态相关约分得到新作用。

\section{一个等式把下一步交给数论}
命题\ref{prop:Ht}只要求证明有限交换群被一批小整数生成。它不再涉及平均的选点过程，已经能直接调用数论工具。下一讲会看到：素数模数有初等证明；一般模数有无条件的大参数结果；在GRH下可以把剩余有限范围接齐。
